\documentclass[10pt,twocolumn]{article}

% ── Packages ──────────────────────────────────────────────────────────────
\usepackage[T1]{fontenc}
\usepackage[utf8]{inputenc}
\usepackage[margin=0.75in, top=0.85in, bottom=0.85in]{geometry}
\usepackage{times}
\usepackage{graphicx}
\usepackage{booktabs}
\usepackage{array}
\usepackage{microtype}
\usepackage{caption}
\usepackage{fancyhdr}
\usepackage{titlesec}
\usepackage{hyperref}
\usepackage{amsmath}
\usepackage{float}

% ── IEEE-like formatting ──────────────────────────────────────────────────
\setlength{\columnsep}{0.25in}
\setlength{\parindent}{1em}
\setlength{\parskip}{0pt}

\titleformat{\section}[block]
  {\normalfont\bfseries\centering\normalsize}{}{0em}
  {\thesection.\enspace\MakeUppercase}
\titlespacing{\section}{0pt}{8pt plus 2pt}{4pt plus 1pt}

\titleformat{\subsection}[block]
  {\normalfont\itshape\normalsize}{}{0em}
  {\Alph{subsection}.\enspace}
\titlespacing{\subsection}{0pt}{6pt plus 1pt}{2pt}

\titleformat{\subsubsection}[runin]
  {\normalfont\itshape\normalsize\bfseries}{}{0em}{}[.]
\titlespacing{\subsubsection}{0pt}{4pt}{0.5em}

% Abstract environment
\newenvironment{ieeeabstract}{%
  \begin{center}\begin{minipage}{0.9\linewidth}
  \small\noindent\textbf{Abstract---}\ignorespaces}
 {\end{minipage}\end{center}\vspace{4pt}}

% Header / footer
\pagestyle{fancy}
\fancyhf{}
\fancyhead[C]{\footnotesize Systems Analysis \& Design --- Workshop No.~1 \& 2, Semester 2026-I}
\fancyfoot[C]{\footnotesize\thepage}
\renewcommand{\headrulewidth}{0.4pt}

% Caption
\captionsetup{font=footnotesize, labelfont=bf, labelsep=period}

% Column type
\newcolumntype{L}[1]{>{\raggedright\arraybackslash\footnotesize}p{#1}}

% References title
\renewcommand{\refname}{\normalfont\bfseries\centering\normalsize REFERENCES}

% =============================================================================
\begin{document}
% =============================================================================

% ── Title block ──────────────────────────────────────────────────────────
\twocolumn[{%
\begin{center}
  {\large\bfseries Designing a Community Security Alert System\\[4pt]
   for a University Campus: Analysis, Architecture, and Lessons Learned}\\[10pt]
  {\normalsize
   Felipe~Jose~Garzon~Herrera,~\textit{20251020132},\quad
   Juan~Esteban~Quintero~Gordillo,~\textit{20251020137},\\[2pt]
   Henry~Samuel~Garrido~Medina,~\textit{20251020125},\quad
   Gabriel~Mateo~Cusba~Marin,~\textit{20251020128}\\[6pt]
   \textit{Systems Analysis \& Design, Semester 2026-I}\\
   \textit{Eng.\ Carlos Andr\'es Sierra, M.Sc.}\\
   \textit{Computer Engineering Program, Universidad Distrital Francisco
   Jos\'e de Caldas, Bogot\'a D.C., Colombia}}\\[10pt]
\end{center}
\begin{ieeeabstract}
University campuses in Latin America face persistent safety challenges
rooted in the absence of real-time, community-driven incident reporting
tools. At Universidad Distrital Francisco Jos\'e de Caldas, systematic data
collection revealed three compounding problems: a strong under-reporting
culture driven by high-friction reporting channels, a concentration of
incidents in specific zones and hours, and a lack of centralized incident
management that produces average security response delays of~11.4~minutes.
This paper presents the Community Security Alert System (CSAS), a
crowd-sourced mobile platform designed to address these gaps through a
microservices-based architecture with a hybrid AI-and-human verification
pipeline and geofenced alert dispatch. The proposed design projects a
reduction in response time to under~5~minutes and an increase in incident
recording completeness from~41\% to above~85\%, establishing a replicable
methodology for campus security systems in similar institutional contexts.
\end{ieeeabstract}
\vspace{6pt}
\noindent\textbf{Index Terms---}campus security, crowd-sourced incident
reporting, microservices, real-time alerts, systems analysis, systems design.
\vspace{10pt}
}]

% =============================================================================
\section{Introduction}
% =============================================================================

Campus security is a problem that most universities acknowledge but few solve
effectively. At Universidad Distrital Francisco Jos\'e de Caldas, a public
university in Bogot\'a D.C., the security team relies on phone calls and
walk-up reports to respond to incidents --- a process that takes an average of
11.4~minutes from the moment an event occurs to the moment the response
team is notified. During that time, a situation can escalate significantly.

Several platforms have attempted to address similar problems in the Bogot\'a
metropolitan area, and their limitations are instructive. The Polic\'ia
Nacional's Vigila App integrates directly with law enforcement but requires
users to place a phone call, making it too slow for fast-moving campus
situations~\cite{sierra2026}. Neighborhood WhatsApp groups achieve high
adoption but operate without any verification mechanism, routinely generating
false alarms and panic~\cite{checkland1981}. SIURE~UD, the university's own
internal system, is familiar to staff but operates exclusively through email
and phone, with no mobile interface or real-time capability~\cite{sierra2026}.
The pattern across these platforms is consistent: each solves part of the
problem but creates a different one in its place.

The theoretical framing for this project comes from two foundational works.
Checkland's soft systems methodology~\cite{checkland1981} helped us treat
campus security not as a purely technical problem but as a sociotechnical
system in which human behavior, institutional culture, and technology interact
in complex ways --- a perspective that proved essential during the analysis
phase. Meadows' work on feedback dynamics~\cite{meadows2008} provided the
vocabulary to understand how a crowd-sourced alert system can either amplify
its own usefulness through reinforcing loops, or undermine it through alert
fatigue if poorly calibrated. On the technical side, the systems engineering
life-cycle framework of ISO/IEC/IEEE 15288:2015~\cite{iso15288} structured our
progression from analysis to design, while Bass, Clements, and
Kazman~\cite{bass2013} and Dragoni et al.~\cite{dragoni2017} informed our
architectural decisions, particularly the choice of a microservices-based
decomposition over a monolithic design. The mixed-methods data collection
strategy was guided by Creswell and Creswell~\cite{creswell2018}, whose
framework for triangulating qualitative and quantitative evidence helped us
assess the validity of findings derived from a relatively small sample.

This paper presents the Community Security Alert System (CSAS), a
crowd-sourced mobile and web platform that addresses the identified gaps. The
paper covers two phases of the project: a systems analysis phase in which we
characterized the problem empirically, and a systems design phase in which we
translated those findings into a technical architecture. Our central argument
is that the design decisions made in each phase are only meaningful because
they are grounded in real data collected from the campus community --- not
assumptions.

% =============================================================================
\section{Methods and Materials}
% =============================================================================

\subsection{Systems Analysis}

To characterize the campus security problem, we designed a data collection
plan that combined six complementary methods simultaneously. The goal was
to avoid relying on any single perspective, since each method captures a
different dimension of the problem. Structured interviews with five security
staff members and two administrative staff members revealed the operational
reality of campus security from the inside: communication delays, lack of
centralized records, and the need for precise incident location data. A
student survey ($n=20$) measured safety perception, incident experience, and
willingness to use a reporting application, and showed that 35\% of respondents
had experienced or witnessed an incident, with all reported events occurring
during afternoon or night hours. Direct observation across four campus zones
--- main entrance, library area, parking lot, and dormitory --- over eight
sessions confirmed that the parking lot at night is the highest-risk zone,
with two suspicious events observed both coinciding with poor lighting and low
foot traffic. Document analysis of institutional security records confirmed
a significant under-reporting gap: official records show far fewer incidents
than the survey and observation data suggest. Benchmarking of three existing
platforms in Bogot\'a extracted transferable design insights while identifying
the gaps that none of them fully close. Finally, weather monitoring over one
week suggested a preliminary correlation between rainy nights and incident
occurrence, motivating a weather-aware component in the design.

The convergence of three independent methods --- interviews, survey, and
observation --- on the same three problems gave us confidence in the findings:
under-reporting due to high friction, night-time and spatial concentration of
incidents, and the absence of a centralized real-time tool. When independent
methods agree, the conclusion is stronger than any single source of evidence
can provide.

Figure~\ref{fig:context} shows the stakeholder context diagram, mapping the
information flows between the six groups that interact with the CSAS: students,
professors and administrative staff, the campus security team, local authorities,
system administrators, and university management.

\begin{figure}[H]
\centering
\includegraphics[width=\columnwidth]{fig_context.png}
\caption{Stakeholder relationships and information flows in the CSAS.
Students are the primary reporters; security staff verify and coordinate;
management consumes analytics.}
\label{fig:context}
\end{figure}

\subsection{System Design}

The design phase translated the three empirical findings into three
corresponding architectural decisions. The under-reporting culture is addressed
by a mobile-first interface that requires no more than three interactions to
submit a report (NFR-06), and by an anonymous reporting option for sensitive
incident types. The spatial and temporal concentration of incidents is addressed
by a geofenced alert dispatcher that applies zone-weighted and time-weighted
priority logic, automatically elevating the urgency of night-hour alerts from
the parking lot and dormitory zones. The absence of a centralized tool is
addressed by a microservices-based, event-driven architecture that integrates
with the existing SIURE~UD institutional system rather than replacing it.

The architecture is organized across five layers. The presentation layer
encompasses the mobile application (React Native) and the administrative web
portal. The application layer is an API gateway (Node.js / Nginx) that handles
authentication, rate limiting, and routing. The service layer contains the five
core microservices --- Incident, Verification, Dispatcher, User, and Analytics
--- which communicate asynchronously through a RabbitMQ message broker to ensure
loose coupling between report ingestion and alert dispatch. The data layer uses
PostgreSQL with the PostGIS spatial extension, which supports the geospatial
queries required for zone-weighted routing and security heatmap generation.
The integration layer exposes authenticated endpoints to SIURE~UD, local
authorities, and a Twilio SMS gateway that provides a fallback delivery channel
for users without smartphone access.

The verification pipeline deserves particular attention because it sits at
the center of the most critical design trade-off in the system. If verification
is too strict, legitimate incidents get delayed. If it is too loose, false alerts
flood the system and users start ignoring everything --- a phenomenon known as
alert fatigue~\cite{meadows2008}. Our solution is a hybrid AI-and-human pipeline:
an AI plausibility scorer performs initial metadata validation and assigns a
confidence score, while a human moderator reviews low-confidence reports and can
override the scorer at any stage. High-severity incident types (assault, medical
emergency) have a priority bypass path that skips AI scoring entirely and goes
directly to the security team. The NFR-08 requirement formalizes the quality
bar: no alert is broadcast as verified without at least two independent community
confirmations or one administrative approval.

\begin{figure}[H]
\centering
\includegraphics[width=0.82\columnwidth]{fig2_processflow.png}
\caption{Main process flow. A submitted incident passes through the
Verification Service; valid incidents trigger geofenced dispatch while
rejected reports are logged for auditing.}
\label{fig:flow}
\end{figure}

Figure~\ref{fig:flow} illustrates this flow. Figure~\ref{fig:arch} shows the
full layered architecture.

\begin{figure}[H]
\centering
\includegraphics[width=\columnwidth]{fig1_architecture.png}
\caption{CSAS layered architecture. The Service Layer hosts the five core
microservices; external integrations connect through the Integration Layer.}
\label{fig:arch}
\end{figure}

One important limitation of our analysis is the survey sample size ($n=20$),
which is too small for statistical generalization. We also did not audit mobile
connectivity coverage across campus zones, which is a prerequisite risk for
the notification delivery architecture. These are acknowledged open points that
the implementation phase must address before deployment.

% =============================================================================
\section{Results and Discussion}
% =============================================================================

\subsection{Gap Analysis}

The clearest result of the analysis phase is a set of measurable operational gaps
between the current state of campus security and the targets the CSAS is designed
to achieve. Table~\ref{tab:gaps} presents these gaps quantitatively. The response
time gap of $-6.4$~minutes is the most critical: reducing from 11.4~minutes to
under 5~minutes would fundamentally change the security team's ability to
intervene during fast-developing incidents.

\begin{table}[H]
\caption{Operational Gap Analysis: Current State vs.\ CSAS Design Targets}
\label{tab:gaps}
\renewcommand{\arraystretch}{1.15}
\centering
\begin{tabular}{L{2.0cm} L{1.5cm} L{1.4cm} L{1.2cm}}
\toprule
\textbf{Dimension} & \textbf{Baseline} & \textbf{Target} & \textbf{Gap} \\
\midrule
Response time      & 11.4\,min & $<$5\,min  & $-$6.4\,min \\
Reporting rate     & 41\%      & $>$85\%    & $+$44\,pp   \\
App adoption       & $\sim$12\%  & $\geq$60\% & $+$48\,pp   \\
Delivery success   & $\sim$87\%  & $>$99.9\%  & $+$12.9\,pp \\
\bottomrule
\end{tabular}
\end{table}

\subsection{Comparison with Existing Platforms}

Table~\ref{tab:bench} compares the CSAS against the three existing platforms
reviewed during benchmarking. The CSAS is the only platform in this comparison
that combines all four capabilities: mobile-first interface, real-time geofenced
dispatch, incident verification, and institutional integration.

\begin{table}[H]
\caption{Feature Comparison: CSAS vs.\ Existing Platforms in Bogot\'a}
\label{tab:bench}
\renewcommand{\arraystretch}{1.15}
\centering
\begin{tabular}{L{1.7cm} L{1.1cm} L{1.5cm} L{1.0cm} L{1.2cm}}
\toprule
\textbf{Platform} & \textbf{Mobile} & \textbf{Real-time} & \textbf{Verified} & \textbf{Institutional} \\
\midrule
Vigila App    & Partial & No  & No  & Yes \\
WhatsApp nets & Yes     & Yes & No  & No  \\
SIURE UD      & No      & No  & Yes & Yes \\
\textbf{CSAS} & \textbf{Yes} & \textbf{Yes} & \textbf{Yes} & \textbf{Yes} \\
\bottomrule
\end{tabular}
\end{table}

\subsection{Key Performance Indicators}

Since no implementation exists at this stage, the KPIs in Table~\ref{tab:kpis}
represent design targets rather than measured outcomes. They are included here
because they formalize the success criteria against which the future system will
be objectively evaluated, and because comparing the baselines with the targets
already constitutes a result: it quantifies exactly how far the current state is
from what is needed, which is useful information in itself.

\begin{table}[H]
\caption{Key Performance Indicators and Measurement Plan}
\label{tab:kpis}
\renewcommand{\arraystretch}{1.15}
\centering
\begin{tabular}{L{2.2cm} L{2.0cm} L{1.5cm} L{0.9cm}}
\toprule
\textbf{KPI} & \textbf{Target} & \textbf{Method} & \textbf{Freq.} \\
\midrule
Alert latency    & $<$30\,s (NFR-01)  & Event logs     & Daily     \\
Delivery rate    & $>$99.9\% (NFR-02) & Push receipts  & Daily     \\
Availability     & 99.9\% (NFR-03)    & Monitoring     & Cont.     \\
Reporting rate   & $>$85\% logged     & DB records     & Monthly   \\
Adoption rate    & $>$60\% students   & MAU analytics  & Monthly   \\
False alarm rate & $\leq$5\%          & Staff review   & Weekly    \\
\bottomrule
\end{tabular}
\end{table}

\subsection{Discussion}

The most important insight from this project is that the technical architecture
is only as strong as the empirical evidence behind it. Every major design decision
--- the three-interaction submission flow, the night-weighted dispatcher, the
RabbitMQ message broker, the hybrid verification pipeline --- was made in direct
response to a specific finding from the analysis phase. If we had started with
the design, we would likely have built something technically correct but practically
useless.

The feedback dynamics identified in the analysis also shape how we think about
deployment. The CSAS depends on a reinforcing loop to become effective: as more
users report incidents, alerts become more useful, which encourages more reporting.
But this loop only activates above a critical adoption threshold of approximately
20--25\% of the campus population. Below that threshold, the system depends almost
entirely on security staff to feed the verification pipeline. This means that
community engagement is not a nice-to-have feature of the CSAS --- it is a
structural prerequisite for the system to work as designed, and it cannot be
achieved by the technical team alone.

Three assumptions in the design remain unvalidated and represent the primary
risks for the implementation phase. First, we assumed that users will report
incidents if the process is easy enough; our survey supports this directionally
(100\% of incident-experienced participants said they would use the app), but
real-world behavior under operational conditions may differ. Second, we assumed
the verification pipeline can operate within a 45-second latency budget; this
depends on AI model performance and infrastructure configuration that have not
been tested. Third, we assumed reliable mobile connectivity across all campus
zones; we did not audit signal coverage, and dead zones could break the
notification delivery guarantees.

% =============================================================================
\section{Conclusions}
% =============================================================================

This paper presented the analysis and design of the Community Security Alert
System, a crowd-sourced platform for campus safety at Universidad Distrital
Francisco Jos\'e de Caldas. The work demonstrates that combining six independent
data collection methods --- even with a small sample --- can produce a robust
and consistent characterization of a sociotechnical problem. When interviews,
surveys, and direct observation converge on the same three problems (under-reporting,
night-time concentration, and absence of a centralized tool), the evidence is
strong enough to make design decisions with confidence. The proposed microservices
architecture addresses each of these problems directly and is grounded in the
empirical evidence at every decision point.

The most important lesson from this project is that systems analysis is not a
formality before the real work begins --- it is the foundation on which the design
depends. The KPI targets, the verification pipeline calibration, the adoption
threshold, and the night-weighted dispatcher all came from data, not intuition.
For future work, we recommend expanding the student survey to at least 200
participants, conducting a campus WiFi coverage audit before finalizing the
notification delivery architecture, and establishing a formal data processing
agreement with SIURE~UD and university legal counsel as prerequisites to
deployment --- because the 60\% adoption target cannot be achieved without
institutional endorsement.

% ── References ───────────────────────────────────────────────────────────
\begin{thebibliography}{9}

\bibitem{sierra2026}
C.~A.~Sierra,
\textit{TeamWork as Computer Engineers --- CS Collaboration Guidelines}.
Bogot\'a: Universidad Distrital Francisco Jos\'e de Caldas, 2026.

\bibitem{checkland1981}
P.~Checkland,
\textit{Systems Thinking, Systems Practice}.
New York: John Wiley \& Sons, 1981.

\bibitem{meadows2008}
D.~H.~Meadows,
\textit{Thinking in Systems: A Primer}.
White River Junction, VT: Chelsea Green Publishing, 2008.

\bibitem{iso15288}
ISO/IEC/IEEE,
\textit{Systems and Software Engineering --- System Life Cycle Processes},
ISO/IEC/IEEE Std.~15288:2015, 2015.

\bibitem{creswell2018}
J.~W.~Creswell and J.~D.~Creswell,
\textit{Research Design: Qualitative, Quantitative, and Mixed Methods Approaches},
5th~ed. Thousand Oaks, CA: SAGE Publications, 2018.

\bibitem{ley1581}
Congreso de Colombia,
\textit{Ley 1581 de 2012 --- R\'egimen General de Protecci\'on de Datos Personales},
2012.

\bibitem{bass2013}
L.~Bass, P.~Clements, and R.~Kazman,
\textit{Software Architecture in Practice}, 3rd~ed.
Boston: Addison-Wesley, 2013.

\bibitem{dragoni2017}
N.~Dragoni \textit{et al.},
``Microservices: Yesterday, Today, and Tomorrow,''
\textit{Present and Ulterior Software Engineering},
pp.~195--216, 2017.

\end{thebibliography}

\end{document}
